\FigureLayoutDeclare{img/2023/new_gaokao_paper_1/q6_solution_fig1.png}{4.5cm}{36bp 37bp 27bp 40bp}

\chapter{2023年新高考I卷}

\section{单选题}

\begin{problem}
已知集合 \(M=\{-2,-1,0,1,2\}\)，\(N=\{x\mid x^2-x-6\ge 0\}\)，则 \(M\cap N=\)（\quad）

\choices
    {\(\{-2,-1,0,1\}\)}
    {\(\{0,1,2\}\)}
    {\(\{-2\}\)}
    {\(\{2\}\)}
\end{problem}

\begin{problem}
已知 \(z=\frac{1-\i}{2+2\i}\)，则 \(z-\overline z=\)（\quad）

\choices
    {\(-\i\)}
    {\(\i\)}
    {\(0\)}
    {\(1\)}
\end{problem}

\begin{problem}
已知向量 \(\bs{a}=(1,1)\)，\(\bs{b}=(1,-1)\)，若 \((\bs{a}+\lambda\bs{b})\perp(\bs{a}+\mu\bs{b})\)，则（\quad）

\choices
    {\(\lambda+\mu=1\)}
    {\(\lambda+\mu=-1\)}
    {\(\lambda\mu=1\)}
    {\(\lambda\mu=-1\)}
\end{problem}

\begin{problem}
设函数 \(f(x)=2^{x(x-a)}\) 在区间 \((0,1)\) 单调递减，则 \(a\) 的取值范围是（\quad）

\choices
    {\((-\infty,-2]\)}
    {\([-2,0)\)}
    {\((0,2]\)}
    {\([2,+\infty)\)}
\end{problem}

\begin{problem}
设椭圆 \(C_1:\frac{x^2}{a^2}+y^2=1\)（\(a>1\)），\(C_2:\frac{x^2}{4}+y^2=1\) 的离心率分别为 \(e_1\)，\(e_2\)，若 \(e_2=\sqrt3\,e_1\)，则 \(a=\)（\quad）

\choices
    {\(\frac{2\sqrt3}{3}\)}
    {\(\sqrt2\)}
    {\(\sqrt3\)}
    {\(\sqrt6\)}
\end{problem}

\begin{problem}
过点 \((0,-2)\) 与圆 \(x^2+y^2-4x-1=0\) 相切的两直线的夹角为 \(\alpha\)，则 \(\sin\alpha=\)（\quad）

\choices
    {\(1\)}
    {\(\frac{\sqrt{15}}{4}\)}
    {\(\frac{\sqrt{10}}{4}\)}
    {\(\frac{\sqrt6}{4}\)}
\end{problem}

\begin{problem}
记 \(S_n\) 为数列 \(\{a_n\}\) 的前 \(n\) 项和. 设甲：\(\{a_n\}\) 为等差数列，乙：\(\left\{\frac{S_n}{n}\right\}\) 为等差数列，则（\quad）

\choices
    {甲是乙的充分条件但不是必要条件}
    {甲是乙的必要条件但不是充分条件}
    {甲是乙的充要条件}
    {甲既不是乙的充分条件也不是乙的必要条件}
\end{problem}

\begin{problem}
已知 \(\sin(\alpha-\beta)=\frac13\)，\(\cos\alpha\sin\beta=\frac16\)，则 \(\cos(2\alpha+2\beta)=\)（\quad）

\choices
    {\(\frac79\)}
    {\(\frac19\)}
    {\(-\frac19\)}
    {\(-\frac79\)}
\end{problem}

\section{多选题}

\begin{problem}
有一组样本数据 \(x_1\)，\(x_2\)，\(\dots\)，\(x_6\)，其中 \(x_1\) 是最小值，\(x_6\) 是最大值，则（\quad）

\choices
    {\(x_2,x_3,x_4,x_5\) 的平均数等于 \(x_1,x_2,\dots,x_6\) 的平均数}
    {\(x_2\)，\(x_3\)，\(x_4\)，\(x_5\) 的中位数等于 \(x_1\)，\(x_2\)，\(\dots\)，\(x_6\) 的中位数}
    {\(x_2\)，\(x_3\)，\(x_4\)，\(x_5\) 的标准差不小于 \(x_1\)，\(x_2\)，\(\dots\)，\(x_6\) 的标准差}
    {\(x_2\)，\(x_3\)，\(x_4\)，\(x_5\) 的极差不大于 \(x_1\)，\(x_2\)，\(\dots\)，\(x_6\) 的极差}
\end{problem}

\begin{problem}
噪声污染问题越来越受到重视. 用声压级来度量噪声的强度，定义声压级 \(L_P=20\lg\frac{p}{p_0}\)，其中常数 \(p_0(p_0>0)\) 是听觉下限阈值，\(p\) 是实际声压. 下表为不同声源的声压级：

\begin{center}
\begin{tabular}{|c|c|c|}
\hline
声源 & 与声源的距离/m & 声压级/dB \\
\hline
燃油汽车 & 10 & 60\textasciitilde90 \\
\hline
混合动力汽车 & 10 & 50\textasciitilde60 \\
\hline
电动汽车 & 10 & 40 \\
\hline
\end{tabular}
\end{center}

已知在距离燃油汽车，混合动力汽车，电动汽车 \(10\text{m}\) 处测得实际声压分别为 \(p_1\)，\(p_2\)，\(p_3\)，则（\quad）

\choices
    {\(p_1\ge p_2\)}
    {\(p_2>10p_3\)}
    {\(p_3=100p_0\)}
    {\(p_1\le 100p_2\)}
\end{problem}

\begin{problem}
已知函数 \(f(x)\) 的定义域为 \(\R\)，且 \(f(xy)=y^2f(x)+x^2f(y)\)，则（\quad）

\choices
    {\(f(0)=0\)}
    {\(f(1)=0\)}
    {\(f(x)\) 是偶函数}
    {\(x=0\) 为 \(f(x)\) 的极小值点}
\end{problem}

\begin{problem}
下列物体中，能够被整体放入棱长为 \(1\)（单位：m）的正方体容器（容器壁厚度忽略不计）内的有（\quad）

\choices
    {直径为 \(0.99\text{m}\) 的球体}
    {所有棱长均为 \(1.4\text{m}\) 的正四面体}
    {底面直径为 \(0.01\text{m}\)，高为 \(1.8\text{m}\) 的圆柱体}
    {底面直径为 \(1.2\text{m}\)，高为 \(0.01\text{m}\) 的圆柱体}
\end{problem}

\section{填空题}

\begin{problem}
某学校开设了 \(4\) 门体育类选修课和 \(4\) 门艺术类选修课. 学生需从这 \(8\) 门课中选 \(2\) 门或 \(3\) 门课，并且每类选修课至少选 \(1\) 门，则不同的选课方案共有\fillinblank{}种. （用数字作答）
\end{problem}

\begin{problem}
在正四棱台 \(ABCD-A_1B_1C_1D_1\) 中，\(AB=2\)，\(A_1B_1=1\)，\(AA_1=\sqrt2\)，则该棱台的体积为\fillinblank{}.
\end{problem}

\begin{problem}
已知函数 \(f(x)=\cos\omega x-1\)（\(\omega>0\)） 在区间 \([0,2\pi]\) 上有且仅有 \(3\) 个零点，则 \(\omega\) 的取值范围是\fillinblank{}.
\end{problem}

\begin{problem}
已知双曲线 \(C:\frac{x^2}{a^2}-\frac{y^2}{b^2}=1\)（\(a>0\)，\(b>0\)） 的左，右焦点分别为 \(F_1\)，\(F_2\). 点 \(A\) 在 \(C\) 上，点 \(B\) 在 \(y\) 轴上，\(F_1A\perp F_1B\)，\(\overrightarrow{F_2A}=-\frac23\overrightarrow{F_2B}\)，则 \(C\) 的离心率为\fillinblank{}.
\end{problem}

\section{解答题}

\begin{problem}
已知在 \(\triangle ABC\) 中，\(A+B=3C\)，\(2\sin(A-C)=\sin B\).
\begin{enumerate}
\item 求 \(\sin A\)；
\item 设 \(AB=5\)，求 \(AB\) 边上的高.
\end{enumerate}
\end{problem}

\begin{problem}
如图，在正四棱柱 \(ABCD-A_1B_1C_1D_1\) 中，\(AB=2\)，\(AA_1=4\). 点 \(A_2\)，\(B_2\)，\(C_2\)，\(D_2\) 分别在棱 \(AA_1\)，\(BB_1\)，\(CC_1\)，\(DD_1\) 上，且 \(AA_2=1\)，\(BB_2=DD_2=2\)，\(CC_2=3\).

\begin{enumerate}
\item 证明：\(B_2C_2\parallel A_2D_2\)；
\item 点 \(P\) 在棱 \(BB_1\) 上，当二面角 \(P-A_2C_2-D_2\) 为 \(150^\circ\) 时，求 \(B_2P\).
\end{enumerate}

\bitmapfigure[max width=0.42\linewidth]{img/2023/new_gaokao_paper_1/q18_fig1.png}
\end{problem}

\begin{problem}
已知函数 \(f(x)=a(\e^x+a)-x\).
\begin{enumerate}
\item 讨论 \(f(x)\) 的单调性；
\item 证明：当 \(a>0\) 时，\(f(x)>2\ln a+\frac32\).
\end{enumerate}
\end{problem}

\begin{problem}
设等差数列 \(\{a_n\}\) 的公差为 \(d\)，且 \(d>1\). 令 \(b_n=\frac{n^2+n}{a_n}\)，记 \(S_n\)，\(T_n\) 分别为数列 \(\{a_n\}\)，\(\{b_n\}\) 的前 \(n\) 项和.
\begin{enumerate}
\item 若 \(3a_2=3a_1+a_3\)，\(S_3+T_3=21\)，求 \(\{a_n\}\) 的通项公式；
\item 若 \(\{b_n\}\) 为等差数列，且 \(S_{99}-T_{99}=99\)，求 \(d\).
\end{enumerate}
\end{problem}

\begin{problem}
甲乙两人投篮，每次由其中一人投篮，规则如下：若命中则此人继续投篮，若未命中则换为对方投篮. 无论之前投篮情况如何，甲每次投篮的命中率均为 \(0.6\)，乙每次投篮的命中率均为 \(0.8\). 由抽签确定第一次投篮的人选，第一次投篮的人是甲，乙的概率各为 \(0.5\).
\begin{enumerate}
\item 求第二次投篮的人是乙的概率；
\item 求第 \(i\) 次投篮的人是甲的概率；
\item 已知：若随机变量 \(X_i\) 服从两点分布，且 \(P(X_i=1)=1-P(X_i=0)=q_i\)，\(\ i=1\)，\(2\)，\(\dots\)，\(n\)，则 \(E\left(\sum_{i=1}^n X_i\right)=\sum_{i=1}^n q_i\).
记前 \(n\) 次（即从第 \(1\) 次到第 \(n\) 次）投篮中甲投篮的次数为 \(Y\)，求 \(E(Y)\).
\end{enumerate}
\end{problem}

\begin{problem}
在直角坐标系 \(xOy\) 中，点 \(P\) 到 \(x\) 轴的距离等于点 \(P\) 到点 \((0,\frac12)\) 的距离，记动点 \(P\) 的轨迹为 \(W\).
\begin{enumerate}
\item 求 \(W\) 的方程；
\item 已知矩形 \(ABCD\) 有三个顶点在 \(W\) 上，证明：矩形 \(ABCD\) 的周长大于 \(3\sqrt3\).
\end{enumerate}
\end{problem}

